\section{Limitations and Future Work}
\label{sec:limitations}

\paragraph{Objectives orthogonal to the declared purpose are invisible by
construction.}
This is the deepest limitation of the approach, and it is structural rather than an
implementation gap. The feature basis is spanned entirely by sub-goals that the
decomposer extracts from the declared purpose $D$. IOR can therefore only detect
divergence along dimensions derived from $D$. A hidden objective that does not
project onto any $D$-derived sub-goal, one that shares no vocabulary with the
declared purpose, produces near-uniform feature scores and reads as low divergence.
The gym's CostMinimiser is detectable only because ``maximise tool-call count''
anti-correlates with natural cost sub-goals such as \texttt{avoids_redundant_calls}
and \texttt{minimises_steps}; an agent whose true objective were genuinely
orthogonal to cost minimisation would evade detection. The blind-critic ablation
(Section~\ref{sec:evaluation}) tests only critic-side bias and does not address this
upstream blindness. Mitigations worth exploring include seeding the basis with
anti-goals and adversarially proposed off-purpose dimensions, but no
purpose-derived basis can be complete, and IOR should be presented as detecting divergence
\emph{relative to a declared purpose}, not as a general objective recovery method.

\paragraph{The bottleneck is judge reliability, not the inference machinery.}
Our own evidence locates the method's centre of gravity in the language-model judge
rather than in the Bayesian IRL. The idealised-critic check recovers the planted
divergence at $r = 0.96$ on the same trajectories where the live judge reaches
$-0.72$ (Section~\ref{subsec:recovery}), so the binding constraint is whether a
language model can reliably score the specific sub-goals at issue. This places IOR
squarely within the study of LLM-as-judge reliability and its documented failure
modes, including verbosity, position, and self-enhancement
bias~\citep{zheng2023judge, llmJudgeSurvey2024}, a literature the method depends on
and which bounds it. Two consequences follow for any use as regulatory evidence.
First, a confident wrong-direction failure ($r = -0.72$) is more dangerous than
abstention: a tool that mis-localises divergence points an auditor at the wrong
behaviour, so calibrated abstention on low-confidence sub-goals is a prerequisite we
do not yet provide. Second, our evidence base is three hand-authored agents with a
single clear success, which is too small and too self-authored to support general
claims; the gym must grow to many agents spanning divergence types, ideally authored
independently of the featuriser, before recovery on behavioural divergences can be
claimed as a property of the method rather than of one example.

\paragraph{The closed loop is designed but not yet demonstrated.}
We describe the inference-probe feedback loop as the architectural contribution, but
we have not yet run an experiment that closes it: the present evaluation compares
targeted against uniform probes in a single shot and does not test whether
re-inferring on probe-generated traces improves the divergence estimate. The loop
also carries a risk we must control for before claiming benefit. Probes are
generated where divergence is already suspected and are designed to make a diverging
agent behave distinctly; feeding the resulting traces back into inference can, in
principle, entrench an initial misestimate rather than correct it. A credible
demonstration therefore needs a control in which the loop is run against an agent
that does \emph{not} diverge, to show that it converges to low divergence rather
than amplifying its own prior. We treat the loop as a proposal pending this
evidence.

\paragraph{Sparse trajectories.}
The Laplace approximation degrades for trajectories shorter than approximately ten
steps, where the Gaussian posterior assumption is least reliable. IOR provides a
Metropolis-Hastings sampler as an alternative, and a future benchmark will compare
it against the Laplace baseline on short gym trajectories.

\paragraph{Stochasticity versus genuine divergence.}
An agent that explores stochastically produces feature vectors that resemble
divergence under BIRL. IOR currently conflates exploration with objective drift; a
principled treatment would require a non-stationary extension of the reward model.

\paragraph{Feature quality.}
The GDJ features are only as good as the judge's ability to decompose and critique.
For highly technical declared purposes or opaque tool semantics, feature quality
may degrade. Calibration on the planted-divergence gym is the current mitigation,
and the structural and ensemble featurisers provide a reproducibility floor and a
robustness check respectively.

\paragraph{Log interoperability.}
IOR's schema is a neutral interchange format, but adapters from provider-specific
log formats (OpenAI, Anthropic, tracing frameworks) into that schema are not yet
implemented. Two schema questions remain: how to represent parallel tool calls
within a single turn, and how to represent text-only reasoning steps that invoke no
tool.

\paragraph{Unvalidated posterior quality and an unspecified taxonomy mapping.}
Two narrower gaps bear on the regulatory-legibility claim. First, the Laplace
approximation underlying the reported uncertainties $\sigma_i$ has now been checked
against the Metropolis-Hastings sampler (Section~\ref{subsec:calibration}) and is
mildly conservative (within $20\%$), but only on the deterministic structural basis;
its calibration on the live-judge path, where scoring is itself non-deterministic,
remains unvalidated. Second, the claim that every finding maps onto an OWASP, MITRE ATLAS,
or NIST node is at present realised by a language-model classifier in the scorer
that assigns a taxonomy node and severity to each probe finding; the mapping from a
divergence dimension to a taxonomy node is not a fixed, audited correspondence but a
model judgement, and a defensible compliance artefact would need that mapping
specified and validated rather than delegated to a model.

\paragraph{Normalisation choice.}
The divergence computation maps reward weights to the probability simplex by a
softmax and takes the total variation distance, having found that an earlier
min-subtraction normalisation forced the lowest-weighted dimension to exactly zero
and inflated its apparent divergence. The choice is not innocuous: the
IRL-versus-mean comparison of Section~\ref{sec:evaluation} shifts materially between
the two projections (Spearman $0.77$ under softmax against $0.91$ under
min-subtraction), so divergence conclusions carry a dependence on the projection
that we have characterised only partially. A KL distance is a further natural
alternative whose effect we have not measured.

\paragraph{Regulatory evidentiary weight.}
Whether ``inferred objective versus intended purpose'' carries evidentiary weight
under EU AI Act Art.~9 and Annex~IV, or serves only as supporting documentation, is
an open legal question. IOR frames its output as supporting evidence rather than a
definitive compliance finding.

\paragraph{Multi-agent settings.}
IOR audits single agents. For multi-agent systems the joint objective may diverge
from any individual agent's declared purpose; the gym includes a cooperating-pair
scenario whose interleaved trajectory is audited as one system, and full joint-IRL
treatment is left to future work.
